% ============================================================================
% APRESENTAÇÃO COMPLEMENTAR: HISTÓRIA E FILOSOFIA DA BIOLOGIA DE SISTEMAS
% Bioinformática para Biologia de Sistemas
% ============================================================================

% ============================================================================
% TEMPLATE BASE PARA APRESENTAÇÕES EM BEAMER
% Bioinformática para Biologia de Sistemas
% ============================================================================

\documentclass[12pt, aspectratio=169]{beamer}

% ============================================================================
% PACOTES ESSENCIAIS
% ============================================================================
\usepackage[utf8]{inputenc}
\usepackage[T1]{fontenc}
\usepackage[brazilian]{babel}
\usepackage{amsmath, amssymb, amsthm}
\usepackage{graphicx}
\usepackage{xcolor}
\usepackage{tikz}
\usetikzlibrary{arrows.meta, shapes, positioning, calc, decorations.pathreplacing, decorations.shapes, decorations.pathmorphing}
\usepackage{listings}
\usepackage{booktabs}
\usepackage{multirow}
\usepackage{hyperref}
\usepackage{chemfig}
\usepackage{chemarr}

% ============================================================================
% CONFIGURAÇÃO DO TEMA
% ============================================================================
\usetheme{Madrid}
\usecolortheme{default}

% Cores personalizadas
\definecolor{azulescuro}{RGB}{0, 51, 102}
\definecolor{azulclaro}{RGB}{51, 153, 255}
\definecolor{verdeacido}{RGB}{102, 204, 0}
\definecolor{laranjacel}{RGB}{255, 153, 51}
\definecolor{roxodna}{RGB}{153, 51, 255}

\setbeamercolor{structure}{fg=azulescuro}
\setbeamercolor{palette primary}{bg=azulescuro,fg=white}
\setbeamercolor{palette secondary}{bg=azulclaro,fg=white}
\setbeamercolor{palette tertiary}{bg=azulescuro,fg=white}
\setbeamercolor{palette quaternary}{bg=azulclaro,fg=white}
\setbeamercolor{block title}{bg=azulescuro,fg=white}
\setbeamercolor{block body}{bg=azulclaro!10,fg=black}
\setbeamercolor{block title alerted}{bg=red!80,fg=white}
\setbeamercolor{block body alerted}{bg=red!10,fg=black}
\setbeamercolor{block title example}{bg=verdeacido!80,fg=white}
\setbeamercolor{block body example}{bg=verdeacido!10,fg=black}

% ============================================================================
% CONFIGURAÇÃO DE CÓDIGO
% ============================================================================
\lstset{
    language=Python,
    basicstyle=\ttfamily\small,
    keywordstyle=\color{azulescuro}\bfseries,
    commentstyle=\color{verdeacido},
    stringstyle=\color{laranjacel},
    numbers=left,
    numberstyle=\tiny\color{gray},
    stepnumber=1,
    numbersep=5pt,
    backgroundcolor=\color{white},
    showspaces=false,
    showstringspaces=false,
    showtabs=false,
    frame=single,
    rulecolor=\color{black},
    tabsize=2,
    captionpos=b,
    breaklines=true,
    breakatwhitespace=false,
    escapeinside={\%*}{*)},
    xleftmargin=10pt,
    xrightmargin=10pt
}

% Definir estilo para R
\lstdefinestyle{Rstyle}{
    language=R,
    basicstyle=\ttfamily\small,
    keywordstyle=\color{azulescuro}\bfseries,
    commentstyle=\color{verdeacido},
    stringstyle=\color{laranjacel}
}

% Definir estilo para MATLAB
\lstdefinestyle{MATLABstyle}{
    language=Matlab,
    basicstyle=\ttfamily\small,
    keywordstyle=\color{azulescuro}\bfseries,
    commentstyle=\color{verdeacido},
    stringstyle=\color{laranjacel}
}

% ============================================================================
% COMANDOS PERSONALIZADOS
% ============================================================================

% Comando para equações destacadas
\newcommand{\destaque}[1]{%
    \begin{alertblock}{}
        \begin{center}
            #1
        \end{center}
    \end{alertblock}
}

% Comando para definições
\newcommand{\definicao}[2]{%
    \begin{block}{Definição: #1}
        #2
    \end{block}
}

% Comando para exemplos
\newcommand{\exemplo}[2]{%
    \begin{exampleblock}{Exemplo: #1}
        #2
    \end{exampleblock}
}

% Comando para observações importantes
\newcommand{\importante}[1]{%
    \begin{alertblock}{Importante!}
        #1
    \end{alertblock}
}

% Comandos matemáticos úteis
\newcommand{\R}{\mathbb{R}}
\newcommand{\N}{\mathbb{N}}
\newcommand{\Z}{\mathbb{Z}}
\newcommand{\dif}{\mathrm{d}}
\newcommand{\parcial}[2]{\frac{\partial #1}{\partial #2}}
\newcommand{\derivada}[2]{\frac{\dif #1}{\dif #2}}
\newcommand{\vect}[1]{\boldsymbol{#1}}

% Comandos biológicos
\newcommand{\gene}[1]{\textit{#1}}
\newcommand{\proteina}[1]{\textsc{#1}}
\newcommand{\especie}[1]{\textit{#1}}

% ============================================================================
% CONFIGURAÇÃO DE RODAPÉ
% ============================================================================
\setbeamertemplate{footline}{
  \leavevmode%
  \hbox{%
  \begin{beamercolorbox}[wd=.33\paperwidth,ht=2.25ex,dp=1ex,center]{author in head/foot}%
    \usebeamerfont{author in head/foot}\insertshortauthor
  \end{beamercolorbox}%
  \begin{beamercolorbox}[wd=.34\paperwidth,ht=2.25ex,dp=1ex,center]{title in head/foot}%
    \usebeamerfont{title in head/foot}\insertshorttitle
  \end{beamercolorbox}%
  \begin{beamercolorbox}[wd=.33\paperwidth,ht=2.25ex,dp=1ex,right]{date in head/foot}%
    \usebeamerfont{date in head/foot}\insertshortdate{}\hspace*{2em}
    \insertframenumber{} / \inserttotalframenumber\hspace*{2ex}
  \end{beamercolorbox}}%
  \vskip0pt%
}

% ============================================================================
% CONFIGURAÇÃO DE NAVEGAÇÃO
% ============================================================================
\setbeamertemplate{navigation symbols}{}

% ============================================================================
% CONFIGURAÇÃO DE ITEMIZE/ENUMERATE
% ============================================================================
\setbeamertemplate{itemize items}[circle]
\setbeamertemplate{enumerate items}[default]

% ============================================================================
% FIM DO TEMPLATE
% ============================================================================


% ============================================================================
% INFORMAÇÕES DO DOCUMENTO
% ============================================================================
\title[História e Filosofia]{Biologia de Sistemas}
\subtitle{Uma Perspectiva Histórica e Filosófica}
\author{Ney Lemke}
\institute{DBF-Unesp}
\date{\today}

\begin{document}

% ============================================================================
% SLIDE DE TÍTULO
% ============================================================================
\begin{frame}
\titlepage
\end{frame}

% ============================================================================
% SUMÁRIO
% ============================================================================
\begin{frame}{Sumário}
\tableofcontents
\end{frame}

% ============================================================================
% SEÇÃO 1: INTRODUÇÃO
% ============================================================================
\section{Introdução}

\begin{frame}{Por que Estudar a História do Pensamento Sistêmico?}
\begin{block}{Motivação}
Compreender as raízes históricas e filosóficas da Biologia de Sistemas nos ajuda a:
\begin{itemize}
    \item Entender os fundamentos conceituais da área
    \item Apreciar a evolução do pensamento científico
    \item Reconhecer padrões recorrentes na ciência
    \item Evitar erros do passado
\end{itemize}
\end{block}

\vspace{0.5cm}

\begin{exampleblock}{Questão Central}
Como passamos de uma visão reducionista para uma abordagem sistêmica da biologia?
\end{exampleblock}
\end{frame}

\begin{frame}{O Desafio do Pensamento Sistêmico (1/2): Reducionismo}
\begin{center}
\begin{alertblock}{Reducionismo}
\textbf{Estratégia:}
\begin{itemize}
    \item Dividir para conquistar
    \item Estudar partes isoladas
    \item Análise detalhada
    \item Causalidade linear
\end{itemize}

\vspace{0.2cm}
\textbf{Sucesso:} Biologia molecular, genética
\end{alertblock}
\end{center}
\end{frame}

\begin{frame}{O Desafio do Pensamento Sistêmico (2/2): Holismo}
\begin{center}
\begin{exampleblock}{Holismo/Sistemas}
\textbf{Estratégia:}
\begin{itemize}
    \item Integrar componentes
    \item Estudar interações
    \item Síntese e emergência
    \item Causalidade circular
\end{itemize}

\vspace{0.2cm}
\textbf{Sucesso:} Redes, dinâmica, regulação
\end{exampleblock}
\end{center}
\end{frame}

\begin{frame}{O Desafio do Pensamento Sistêmico (Conclusão)}
\importante{
Ambas as abordagens são necessárias e complementares!
}
\end{frame}

% ============================================================================
% SEÇÃO 2: RAÍZES FILOSÓFICAS ANTIGAS
% ============================================================================
\section{Raízes Filosóficas Antigas}

\begin{frame}{Aristóteles (384-322 a.C.) (1/2): Vida e Obra}
\begin{columns}
\column{0.3\textwidth}
\begin{center}
\textbf{Aristóteles}\\
\vspace{0.2cm}
{\small Filósofo grego}\\
{\small Pai da biologia}
\end{center}

\column{0.7\textwidth}
\begin{block}{Contribuições Fundamentais}
\begin{itemize}
    \item \textbf{Holismo:} ``O todo é maior que a soma das partes''
    \item \textbf{Teleologia:} Conceito de finalidade (\textit{telos})
    \item \textbf{Classificação:} Primeiro sistema de taxonomia
    \item \textbf{Observação:} Ênfase no estudo empírico
\end{itemize}
\end{block}
\end{columns}
\end{frame}

\begin{frame}{Aristóteles (384-322 a.C.) (2/2): O Todo e as Partes}
\begin{exampleblock}{Citação Clássica}
\textit{``O todo é mais do que a mera soma de suas partes.''}\\
--- Aristóteles, \textit{Metafísica}
\end{exampleblock}
\end{frame}

\begin{frame}{Visão Aristotélica da Natureza}
\begin{block}{Quatro Causas}
Aristóteles propôs que para entender qualquer fenômeno, devemos considerar:
\begin{enumerate}
    \item \textbf{Causa Material:} Do que é feito?
    \item \textbf{Causa Formal:} Qual é sua forma/estrutura?
    \item \textbf{Causa Eficiente:} O que o produziu?
    \item \textbf{Causa Final:} Qual é seu propósito?
\end{enumerate}
\end{block}

\vspace{0.3cm}

\begin{exampleblock}{Relevância Moderna}
A biologia de sistemas moderna recupera aspectos dessa visão:
\begin{itemize}
    \item Estrutura (causa formal) $\rightarrow$ Topologia de redes
    \item Propósito (causa final) $\rightarrow$ Função biológica
    \item Integração de múltiplas perspectivas
\end{itemize}
\end{exampleblock}
\end{frame}

\begin{frame}{Outras Tradições Filosóficas (1/2): Filosofia Oriental}
\begin{block}{Taoísmo (China)}
\begin{itemize}
    \item Yin e Yang (dualidade)
    \item Harmonia e equilíbrio
    \item Interconexão universal
    \item Pensamento holístico
\end{itemize}

\vspace{0.3cm}

\textbf{Conceito de Qi:}\\
Energia vital que flui através de sistemas vivos
\end{block}
\end{frame}

\begin{frame}{Outras Tradições Filosóficas (2/2): Filosofia Ocidental}
\begin{block}{Platão}
\begin{itemize}
    \item Teoria das Formas
    \item Idealismo
    \item Padrões universais
\end{itemize}
\end{block}

\vspace{0.3cm}

\begin{block}{Estoicismo}
\begin{itemize}
    \item Cosmos como organismo
    \item Interconexão causal
    \item Logos (razão universal)
\end{itemize}
\end{block}
\end{frame}

% ============================================================================
% SEÇÃO 3: A REVOLUÇÃO CIENTÍFICA E O REDUCIONISMO
% ============================================================================
\section{A Revolução Científica e o Reducionismo}

\begin{frame}{René Descartes (1596-1650): O Pai da Modernidade}
\begin{columns}
\column{0.3\textwidth}
\begin{center}
\textbf{René Descartes}\\
\vspace{0.2cm}
{\small Filósofo e matemático}\\
{\small ``Penso, logo existo''}
\end{center}

\column{0.7\textwidth}
\begin{block}{O Ponto de Virada}
Descartes estabelece as bases do método científico moderno, fundamentado na razão e na análise:
\begin{itemize}
    \item Ruptura com a tradição escolástica
    \item Busca por certezas indubitáveis
    \item Aplicação do rigor matemático à natureza
\end{itemize}
\end{block}
\end{columns}
\end{frame}

\begin{frame}{O Método de Análise: Dividir para Compreender}
\begin{exampleblock}{O Segundo Preceito do Método}
\textit{``Dividir cada uma das dificuldades que eu examinasse em tantas parcelas quanto possível e quanto fosse necessário para melhor as resolver.''}
\begin{flushright}
--- Descartes, \textit{Discurso sobre o Método} (1637)
\end{flushright}
\end{exampleblock}

\vspace{0.3cm}

\begin{block}{O Nascimento do Reducionismo}
Esta diretriz define a essência do reducionismo:
\begin{itemize}
    \item A compreensão do sistema vem da compreensão de suas partes
    \item Problemas complexos são decompostos em elementos simples
    \item A síntese é vista como a reversão da análise
\end{itemize}
\end{block}
\end{frame}

\begin{frame}{O Mundo como Máquina: A Filosofia Mecanicista}
\begin{block}{Mecanicismo Universal}
Para Descartes, a natureza (incluindo organismos vivos) funciona como um relógio:
\begin{itemize}
    \item \textbf{Res Extensa:} A matéria é pura extensão e movimento
    \item \textbf{Automatismo Animal:} Animais como autômatos biológicos complexos
    \item \textbf{Causalidade Linear:} Ação por contato e movimento mecânico
\end{itemize}
\end{block}

\vspace{0.3cm}

\importante{
Essa visão eliminou as ``causas finais'' de Aristóteles da física e da biologia experimental
}
\end{frame}

\begin{frame}{Dualismo e o Legado Reducionista}
\begin{block}{Dualismo Cartesiano}
Separação entre \textit{Res Cogitans} (mente) e \textit{Res Extensa} (corpo):
\begin{itemize}
    \item O corpo humano é uma máquina passível de estudo anatômico e físico
    \item A biologia torna-se o estudo dos mecanismos materiais da vida
\end{itemize}
\end{block}

\vspace{0.3cm}

\begin{exampleblock}{Impacto de Longo Prazo}
\begin{itemize}
    \item \textbf{Sucesso:} Permitiu o avanço da fisiologia e medicina experimental
    \item \textbf{Limitação:} Dificuldade em lidar com interações complexas e emergência
    \item \textbf{Padrão:} Dominou o pensamento científico por mais de 300 anos
\end{itemize}
\end{exampleblock}
\end{frame}

% ============================================================================
% SEÇÃO 4: BERTALANFFY E A TEORIA GERAL DOS SISTEMAS
% ============================================================================
\section{Bertalanffy e a Teoria Geral dos Sistemas}

\begin{frame}{Ludwig von Bertalanffy (1901-1972)}
\begin{columns}
\column{0.3\textwidth}
\begin{center}
\textbf{L. von Bertalanffy}\\
\vspace{0.2cm}
{\small Biólogo austríaco}\\
{\small Pai da TGS}
\end{center}

\column{0.7\textwidth}
\begin{block}{Obra Fundamental}
\textbf{``General System Theory''} (1968)
\end{block}
\end{columns}
\end{frame}

\begin{frame}{Ludwig von Bertalanffy: Legado}
\importante{
Bertalanffy foi o primeiro a formalizar matematicamente a teoria de sistemas aplicada à biologia
}
\end{frame}

\begin{frame}{Motivação da Teoria Geral dos Sistemas}
\begin{block}{Por que uma teoria geral?}
\begin{itemize}
    \item Crítica ao reducionismo mecanicista
    \item Necessidade de princípios unificadores
    \item Aplicação interdisciplinar
    \item Foco em organismos vivos
\end{itemize}
\end{block}
\end{frame}

\begin{frame}{Conceitos Fundamentais da TGS (1/2)}
\begin{block}{1. Sistema Aberto vs. Fechado}
\begin{itemize}
    \item \textbf{Sistema Fechado:} Sem troca com ambiente (termodinâmica clássica)
    \item \textbf{Sistema Aberto:} Troca matéria e energia (organismos vivos)
\end{itemize}
\end{block}

\begin{block}{2. Estado Estacionário}
\begin{itemize}
    \item Equilíbrio dinâmico (não termodinâmico)
    \item Fluxo contínuo de matéria e energia
    \item Homeostase biológica
\end{itemize}
\end{block}
\end{frame}

\begin{frame}{Conceitos Fundamentais da TGS (2/2)}
\begin{block}{3. Equifinalidade}
\begin{itemize}
    \item Mesmo estado final pode ser alcançado por diferentes caminhos
    \item Independência de condições iniciais
    \item Robustez biológica
\end{itemize}
\end{block}
\end{frame}

\begin{frame}{Mais Conceitos da TGS}
\begin{block}{4. Isomorfismo}
Princípios estruturais semelhantes em diferentes sistemas:
\begin{itemize}
    \item Feedback em circuitos eletrônicos $\leftrightarrow$ Regulação gênica
    \item Redes sociais $\leftrightarrow$ Redes metabólicas
    \item Hierarquia em organizações $\leftrightarrow$ Níveis biológicos
\end{itemize}
\end{block}

\vspace{0.3cm}

\begin{block}{5. Organização Hierárquica}
\begin{itemize}
    \item Sistemas compostos de subsistemas
    \item Propriedades emergentes em cada nível
    \item Moléculas $\rightarrow$ Células $\rightarrow$ Tecidos $\rightarrow$ Organismos
\end{itemize}
\end{block}
\end{frame}

\begin{frame}{Crítica ao Reducionismo (1/2): O Argumento}
\begin{exampleblock}{Argumento de Bertalanffy}
\textit{``A concepção mecanicista encontrou seu triunfo na física moderna. Mas precisamente por causa desse desenvolvimento, tornou-se evidente que a física não pode servir como modelo para todas as ciências.''}
\end{exampleblock}
\end{frame}

\begin{frame}{Crítica ao Reducionismo (2/2): Limitações}
\begin{block}{Limitações do Reducionismo}
\begin{itemize}
    \item Perde-se informação sobre interações
    \item Propriedades emergentes não são previsíveis
    \item Contexto é fundamental
    \item Causalidade circular em sistemas vivos
\end{itemize}
\end{block}

\vspace{0.3cm}

\importante{
Bertalanffy não rejeitava o reducionismo, mas argumentava pela necessidade de complementá-lo com uma visão sistêmica
}
\end{frame}

% ============================================================================
% SEÇÃO 4: DESENVOLVIMENTO DO PENSAMENTO SISTÊMICO (SÉC. XX)
% ============================================================================
\section{Desenvolvimento do Pensamento Sistêmico}

\begin{frame}{Cibernética: Norbert Wiener (1/2)}
\begin{block}{Cibernética (1948)}
Ciência do controle e comunicação em animais e máquinas

\vspace{0.3cm}

\textbf{Conceitos-chave:}
\begin{itemize}
    \item \textbf{Feedback:} Retroalimentação positiva e negativa
    \item \textbf{Controle:} Regulação automática
    \item \textbf{Informação:} Padrões que reduzem incerteza
    \item \textbf{Homeostase:} Manutenção de estados
\end{itemize}
\end{block}
\end{frame}

\begin{frame}{Cibernética: Aplicações Biológicas (2/2)}
\begin{exampleblock}{Aplicações Biológicas}
\begin{itemize}
    \item Regulação da temperatura corporal
    \item Controle de glicose no sangue
    \item Regulação gênica (operons)
    \item Sistemas nervoso e endócrino
\end{itemize}
\end{exampleblock}
\end{frame}

\begin{frame}{Teoria da Informação: Claude Shannon (1/2)}
\begin{block}{``A Mathematical Theory of Communication'' (1948)}
\textbf{Contribuições:}
\begin{itemize}
    \item Quantificação da informação (bits)
    \item Entropia informacional
    \item Capacidade de canal
    \item Codificação e redundância
\end{itemize}
\end{block}
\end{frame}

\begin{frame}{Teoria da Informação: Impacto (2/2)}
\begin{exampleblock}{Impacto na Biologia}
\begin{itemize}
    \item \textbf{Código genético:} Informação em DNA
    \item \textbf{Entropia:} Medida de desordem/informação
    \item \textbf{Redundância:} Códons degenerados
    \item \textbf{Comunicação celular:} Sinalização
\end{itemize}
\end{exampleblock}

\vspace{0.3cm}

\importante{
A biologia molecular moderna é fundamentalmente uma ciência da informação
}
\end{frame}

\begin{frame}{Teoria da Complexidade (1/2): Prigogine}
\begin{block}{Ilya Prigogine (1917-2003)}
\textbf{Nobel de Química (1977)}

\vspace{0.2cm}

\textbf{Estruturas Dissipativas:}
\begin{itemize}
    \item Auto-organização
    \item Longe do equilíbrio
    \item Ordem a partir do caos
    \item Irreversibilidade
\end{itemize}

\vspace{0.2cm}
\textbf{Exemplo:} Reação de Belousov-Zhabotinsky
\end{block}
\end{frame}

\begin{frame}{Teoria da Complexidade (2/2): Kauffman}
\begin{block}{Stuart Kauffman}
\textbf{Biólogo teórico}

\vspace{0.2cm}

\textbf{Conceitos:}
\begin{itemize}
    \item Auto-organização
    \item Redes booleanas
    \item ``Order for free''
    \item Paisagens adaptativas
\end{itemize}

\vspace{0.2cm}
\textbf{Obra:} \textit{The Origins of Order} (1993)
\end{block}
\end{frame}

\begin{frame}{Complexidade e Biologia: Implicação}
\begin{exampleblock}{Implicação}
Ordem e complexidade podem emergir espontaneamente em sistemas biológicos
\end{exampleblock}
\end{frame}

\begin{frame}{Sistemas Adaptativos Complexos (1/2): Características}
\begin{block}{Características}
\begin{enumerate}
    \item \textbf{Agentes:} Múltiplos componentes interagindo
    \item \textbf{Não-linearidade:} Efeitos desproporcionais
    \item \textbf{Emergência:} Propriedades do sistema $\neq$ soma das partes
    \item \textbf{Adaptação:} Aprendizado e evolução
    \item \textbf{Auto-organização:} Padrões sem controle central
\end{enumerate}
\end{block}
\end{frame}

\begin{frame}{Sistemas Adaptativos Complexos (2/2): Exemplos}
\begin{exampleblock}{Exemplos Biológicos}
\begin{itemize}
    \item Sistema imunológico
    \item Redes neurais
    \item Ecossistemas
    \item Evolução
    \item Desenvolvimento embrionário
\end{itemize}
\end{exampleblock}
\end{frame}

% ============================================================================
% SEÇÃO 5: BIOLOGIA DE SISTEMAS MODERNA
% ============================================================================
\section{Biologia de Sistemas Moderna}

\begin{frame}{Definição Moderna: Hiroaki Kitano}
\begin{exampleblock}{Kitano (2002) - Science}
\textit{``Systems biology is an approach to understanding biological systems that emphasizes the study of the interactions and relationships between the components of a system, rather than the components themselves.''}
\end{exampleblock}

\vspace{0.5cm}

\begin{block}{Quatro Pilares da Biologia de Sistemas}
\begin{enumerate}
    \item \textbf{Estrutura do sistema:} Topologia de redes
    \item \textbf{Dinâmica do sistema:} Comportamento temporal
    \item \textbf{Método de controle:} Regulação e feedback
    \item \textbf{Método de design:} Princípios de organização
\end{enumerate}
\end{block}
\end{frame}

\begin{frame}{Pioneiros da Biologia de Sistemas (1/2)}
\begin{block}{Denis Noble}
\textbf{Fisiologista britânico}

\vspace{0.2cm}

\textbf{Contribuições:}
\begin{itemize}
    \item Modelagem cardíaca
    \item Fisiologia de sistemas
    \item Crítica ao gene-centrismo
\end{itemize}

\vspace{0.2cm}
\textbf{Obra:} \textit{The Music of Life} (2006)
\end{block}
\end{frame}

\begin{frame}{Pioneiros da Biologia de Sistemas (2/2)}
\begin{block}{Uri Alon}
\textbf{Biofísico israelense}

\vspace{0.2cm}

\textbf{Contribuições:}
\begin{itemize}
    \item Network motifs
    \item Princípios de design
    \item Robustez
\end{itemize}

\vspace{0.2cm}
\textbf{Obra:} \textit{An Introduction to Systems Biology} (2006)
\end{block}
\end{frame}

\begin{frame}{Conceitos Chave: Motifs}
\begin{exampleblock}{Network Motifs}
Padrões recorrentes em redes biológicas: feedforward loops, feedback loops, etc.
\end{exampleblock}
\end{frame}

\begin{frame}{Mais Autores Modernos (1/2): Kirschner \& Gerhart}
\begin{block}{Marc Kirschner \& John Gerhart}
\textbf{Conceitos:}
\begin{itemize}
    \item \textbf{Robustez:} Resistência a perturbações
    \item \textbf{Evolvabilidade:} Capacidade de evolução
    \item \textbf{Modularidade:} Organização em módulos
\end{itemize}

\vspace{0.2cm}

\textbf{Obra:} \textit{The Plausibility of Life} (2005)
\end{block}
\end{frame}

\begin{frame}{Mais Autores Modernos (2/2): Outros}
\begin{block}{Outros Contribuidores Importantes}
\begin{itemize}
    \item \textbf{Leroy Hood:} Biologia de sistemas e medicina de precisão
    \item \textbf{Albert-László Barabási:} Teoria de redes complexas
    \item \textbf{James Collins:} Biologia sintética
    \item \textbf{Bernhard Palsson:} Modelagem metabólica
\end{itemize}
\end{block}
\end{frame}

\begin{frame}{Abordagens Contemporâneas (1/2)}
\begin{columns}
\column{0.5\textwidth}
\begin{block}{Biologia de Redes}
\begin{itemize}
    \item Redes metabólicas
    \item Redes de regulação gênica
    \item Redes de interação proteína-proteína
    \item Análise topológica
\end{itemize}
\end{block}

\column{0.5\textwidth}
\begin{block}{Modelagem Matemática}
\begin{itemize}
    \item ODEs (equações diferenciais)
    \item Modelos estocásticos
    \item Constraint-based models
    \item Agent-based models
\end{itemize}
\end{block}
\end{columns}
\end{frame}

\begin{frame}{Abordagens Contemporâneas (2/2)}
\begin{columns}
\column{0.5\textwidth}
\begin{block}{Biologia Sintética}
\begin{itemize}
    \item Design de circuitos
    \item Engenharia de sistemas
    \item Standardização
    \item BioBricks
\end{itemize}
\end{block}

\column{0.5\textwidth}
\begin{block}{Medicina de Sistemas}
\begin{itemize}
    \item Medicina personalizada
    \item Biomarcadores
    \item Redes de doenças
    \item Terapias sistêmicas
\end{itemize}
\end{block}
\end{columns}
\end{frame}

% ============================================================================
% SEÇÃO 6: O DESAFIO DO PENSAMENTO SISTÊMICO
% ============================================================================
\section{O Desafio do Pensamento Sistêmico}

\begin{frame}{Reducionismo vs. Holismo: Uma Síntese (1/2)}
\begin{block}{Reducionismo}
\textbf{Vantagens:}
\begin{itemize}
    \item Precisão e rigor
    \item Controle experimental
    \item Mecanismos detalhados
\end{itemize}

\textbf{Limitações:}
\begin{itemize}
    \item Perde contexto
    \item Ignora emergência
    \item Fragmentação do conhecimento
\end{itemize}
\end{block}
\end{frame}

\begin{frame}{Reducionismo vs. Holismo: Uma Síntese (2/2)}
\begin{block}{Holismo/Sistemas}
\textbf{Vantagens:}
\begin{itemize}
    \item Visão integrada
    \item Captura emergência
    \item Contexto e função
\end{itemize}

\textbf{Limitações:}
\begin{itemize}
    \item Complexidade
    \item Dificuldade experimental
    \item Risco de superficialidade
\end{itemize}
\end{block}
\end{frame}

\begin{frame}{Propriedades Emergentes (1/2): Definição}
\definicao{Emergência}{
Propriedade ou comportamento de um sistema que não pode ser deduzido apenas do conhecimento de seus componentes individuais.
}

\vspace{0.5cm}

\importante{
``More is different'' --- Philip Anderson (Nobel de Física, 1977)
}
\end{frame}

\begin{frame}{Propriedades Emergentes (2/2): Exemplos}
\begin{exampleblock}{Exemplos em Biologia}
\begin{itemize}
    \item \textbf{Consciência:} Emerge de neurônios individuais
    \item \textbf{Vida:} Emerge de moléculas não-vivas
    \item \textbf{Ritmos circadianos:} Emerge de loops de feedback
    \item \textbf{Morfogênese:} Padrões emergem de regras locais
    \item \textbf{Comportamento social:} Emerge de indivíduos
\end{itemize}
\end{exampleblock}
\end{frame}

\begin{frame}{Mudança de Paradigma (1/2): Tradicional}
\begin{center}
\begin{alertblock}{Paradigma Tradicional}
\begin{itemize}
    \item Análise
    \item Linearidade
    \item Determinismo
    \item Equilíbrio
    \item Isolamento
    \item Estático
\end{itemize}
\end{alertblock}
\end{center}
\end{frame}

\begin{frame}{Mudança de Paradigma (2/2): Sistêmico}
\begin{center}
\begin{exampleblock}{Paradigma Sistêmico}
\begin{itemize}
    \item Síntese
    \item Não-linearidade
    \item Estocasticidade
    \item Não-equilíbrio
    \item Interação
    \item Dinâmico
\end{itemize}
\end{exampleblock}
\end{center}
\end{frame}

\begin{frame}{Implicações da Mudança de Paradigma}
\begin{block}{Implicações}
\begin{itemize}
    \item Necessidade de novas ferramentas matemáticas
    \item Integração de dados multi-ômicos
    \item Colaboração interdisciplinar
    \item Mudança na educação científica
\end{itemize}
\end{block}
\end{frame}

\begin{frame}{Desafios Atuais}
\begin{block}{Desafios Técnicos}
\begin{itemize}
    \item \textbf{Dados:} Integração de big data multi-escala
    \item \textbf{Modelos:} Validação e parametrização
    \item \textbf{Computação:} Simulações de larga escala
    \item \textbf{Experimentação:} Perturbações sistêmicas
\end{itemize}
\end{block}

\vspace{0.3cm}

\begin{block}{Desafios Conceituais}
\begin{itemize}
    \item \textbf{Causalidade:} Circular vs. linear
    \item \textbf{Previsibilidade:} Limites em sistemas complexos
    \item \textbf{Explicação:} Mecanismo vs. princípio
    \item \textbf{Redução:} Quando é apropriada?
\end{itemize}
\end{block}
\end{frame}

% ============================================================================
% SEÇÃO 7: CONCLUSÃO
% ============================================================================
\section{Conclusão}

\begin{frame}{Síntese Histórica (1/2): Evolução}
\begin{block}{Evolução do Pensamento}
\begin{enumerate}
    \item \textbf{Antiguidade:} Holismo intuitivo (Aristóteles)
    \item \textbf{Revolução Científica:} Triunfo do reducionismo
    \item \textbf{Século XX:} Reconhecimento das limitações
    \item \textbf{Bertalanffy:} Formalização da teoria de sistemas
\end{enumerate}
\end{block}
\end{frame}

\begin{frame}{Síntese Histórica (2/2): Consolidação}
\begin{block}{Evolução do Pensamento (Continuação)}
\begin{enumerate}
    \setcounter{enumi}{4}
    \item \textbf{Cibernética:} Feedback e controle
    \item \textbf{Complexidade:} Auto-organização e emergência
    \item \textbf{Século XXI:} Biologia de sistemas como disciplina
\end{enumerate}
\end{block}

\vspace{0.3cm}

\importante{
A história mostra um movimento pendular entre reducionismo e holismo, com síntese gradual
}
\end{frame}

\begin{frame}{Relevância para a Bioinformática (1/2)}
\begin{block}{Bioinformática como Ponte}
A bioinformática moderna é essencial para a biologia de sistemas:

\vspace{0.3cm}
\begin{itemize}
    \item \textbf{Dados:} Gerenciamento de dados multi-ômicos
    \item \textbf{Análise:} Ferramentas para redes e modelagem
    \item \textbf{Visualização:} Representação de sistemas complexos
    \item \textbf{Simulação:} Modelos computacionais
    \item \textbf{Integração:} Conectar níveis biológicos
\end{itemize}
\end{block}
\end{frame}

\begin{frame}{Relevância para a Bioinformática (2/2): Habilidades}
\begin{exampleblock}{Habilidades Necessárias}
\begin{itemize}
    \item Pensamento sistêmico
    \item Modelagem matemática
    \item Programação e algoritmos
    \item Estatística e machine learning
    \item Conhecimento biológico profundo
\end{itemize}
\end{exampleblock}
\end{frame}

\begin{frame}{Futuro da Biologia de Sistemas (1/2)}
\begin{block}{Tendências Emergentes}
\begin{itemize}
    \item \textbf{IA e Machine Learning:} Descoberta de padrões
    \item \textbf{Single-cell omics:} Resolução celular
    \item \textbf{Spatial omics:} Contexto espacial
    \item \textbf{Digital twins:} Modelos personalizados
    \item \textbf{Whole-cell models:} Simulação completa de células
\end{itemize}
\end{block}
\end{frame}

\begin{frame}{Futuro da Biologia de Sistemas (2/2)}
\begin{exampleblock}{Aplicações Futuras}
\begin{itemize}
    \item Medicina de precisão
    \item Terapias sistêmicas
    \item Biologia sintética avançada
    \item Bioengenharia
    \item Compreensão da origem da vida
\end{itemize}
\end{exampleblock}
\end{frame}

\begin{frame}{Mensagem Final}
\begin{exampleblock}{Citação de Bertalanffy}
\textit{``A concepção organísmica é básica para a biologia moderna. Ela é necessária para lidar com problemas teóricos de toda ordem.''}
\end{exampleblock}

\vspace{0.5cm}

\begin{block}{Reflexão}
\begin{itemize}
    \item O pensamento sistêmico não é novo, mas sua formalização é recente
    \item Reducionismo e holismo são complementares, não excludentes
    \item A biologia de sistemas representa uma síntese madura
    \item O futuro da biologia é necessariamente sistêmico e integrativo
\end{itemize}
\end{block}

\vspace{0.3cm}

\begin{center}
\Large{\textbf{Obrigado!}}
\end{center}
\end{frame}

\end{document}
